% 张三慧《大学物理学》第4版，《力学、热学》 第4章《功和能》
% 原书PDF物理页151--155，原书印刷页135--139，习题4.1--4.28。
% 图4.49--4.60的几何依赖用文字标记，矢量重绘待独立校验。
% 外部教材候选区：未作出版授权与独立物理专家审核。
\section*{第4章\quad 功和能}
\subsection*{章末习题（4.1--4.28）}

\paragraph{习题 4.1}
电梯由一个起重间与一个配重组成。
它们分别系在一根绕过定滑轮的钢缆的两端（图4.49）。
起重间（包括负载）质量为 \(M=1200\,\mathrm{kg}\)，
配重质量为 \(m=1000\,\mathrm{kg}\)。
此电梯由和定滑轮同轴的电动机所驱动。
假定起重间由低层从静止开始加速上升，
加速度 \(a=1.5\,\mathrm{m/s^2}\)。
（1）这时滑轮两侧钢缆的拉力各为多少？
（2）加速时间 \(t=1.0\,\mathrm s\)，在此时间内由电动机所做功是多少？
（忽略滑轮与钢缆质量）
（3）在加速 \(t=1.0\,\mathrm s\) 以后，
起重间匀速上升。求它再上升 \(\Delta h=10\,\mathrm m\)
的过程中，电动机又做了多少功？

\paragraph{习题 4.2}
一匹马拉着雪橇沿着冰雪覆盖的圆弧形路面极缓慢地匀速移动。
设圆弧路面的半径为 \(R\)（图4.50），
马对雪橇的拉力总是平行于路面，
雪橇质量为 \(m\)，与路面间的滑动摩擦系数为 \(\mu_k\)。
当把雪橇由底端拉上 \(45^\circ\) 圆弧时，
马对雪橇做功多少？重力和摩擦力各做功多少？
\textit{图示：}圆弧中心位于底端的正上方，底端切线水平，
终点半径偏离竖直向下方向 \(45^\circ\)。

\paragraph{习题 4.3}
2001年9月11日美国纽约世贸中心双子塔遭恐怖分子劫持的飞机撞击而被撞毁
（图4.51）。据美国官方发表的数据，
撞击南楼的飞机是波音767客机，质量为 \(132\,\mathrm t\)，
速度为 \(942\,\mathrm{km/h}\)。
求该客机的动能，这一能量相当于多少TNT炸药的爆炸能量？
\textit{本书表4.1在教材PDF第131页给出1 kg TNT爆炸约\(4.6\times10^6\,\mathrm J\)。}

\paragraph{习题 4.4}
矿砂由料槽均匀落在水平运动的传送带上，
落砂流量 \(q=50\,\mathrm{kg/s}\)，
传送带匀速移动，速率为 \(v=1.5\,\mathrm{m/s}\)。
求电动机拖动皮带的功率，
这一功率是否等于单位时间内落砂获得的动能？为什么？

\paragraph{习题 4.5}
如图4.52所示，A和B两物体质量 \(m_A=m_B\)，
物体B与桌面间滑动摩擦系数 \(\mu_k=0.20\)，
滑轮摩擦不计。
试利用功能概念求物体A自静止落下 \(h=1.0\,\mathrm m\)
时的速度。

\paragraph{习题 4.6}
如图4.53所示，一木块 \(M\) 静止在光滑地平面上。
一子弹 \(m\) 沿水平方向以速度 \(v\) 射入木块内前进一段距离 \(s'\)
而停在木块内，使木块移动了 \(s_1\) 的距离。
（1）相对于地面参考系，在这一过程中子弹和木块的动能变化各是多少？
子弹和木块间的摩擦力对子弹和木块各做了多少功？
（2）证明子弹和木块的总机械能的增量等于
一对摩擦力之一沿相对位移 \(s'\) 做的功。

\paragraph{习题 4.7}
参考系 \(S'\) 相对于参考系 \(S\) 以速度 \(u\) 作匀速运动。
试用伽利略变换和牛顿定律证明：
如果在参考系 \(S\) 中，动能定理式（4.9）成立，
则在参考系 \(S'\) 中形式完全相同的动能定理也成立，
即有
\[
 A'_{AB}=\frac12mv_B'^2-\frac12mv_A'^2.
\]
注意：相对于两个参考系的功以及动能并不相等。

\paragraph{习题 4.8}
一竖直悬挂的弹簧（劲度系数为 \(k\)）下端挂一物体，
平衡时弹簧已有一伸长。
若以物体的平衡位置为竖直 \(y\) 轴的原点，
相应位置作为弹性势能和重力势能的零点，
试证：当物体的位置坐标为 \(y\) 时，
弹性势能和重力势能之和为
\[
 U(y)=\frac12ky^2.
 \]

\paragraph{习题 4.9}
一轻质弹簧原长 \(l_0\)，劲度系数为 \(k\)，
上端固定，下端挂一质量为 \(m\) 的物体，
先用手托住，使弹簧保持原长；
松手后，物体下降的最低位置时弹簧最大伸长量是多少？
物体经过平衡位置时的速率多大？

\paragraph{习题 4.10}
图4.54表示质量为 \(72\,\mathrm{kg}\) 的人跳蹦极。
弹性蹦极带原长 \(20\,\mathrm m\)，
劲度系数为 \(60\,\mathrm{N/m}\)，忽略空气阻力。
（1）此人自跳台跳出后，落下多高时速度最大？
此最大速度是多大？
（2）已知跳台高于下面的水面 \(60\,\mathrm m\)，
此人跳下后会不会触到水面？

\paragraph{习题 4.11}
如原书图4.55，一轻质弹簧劲度系数为 \(k\)，
两端各固定一质量均为 \(M\) 的物块A和B，
放在水平光滑桌面上静止。
今有一质量为 \(m\) 的子弹沿弹簧轴线方向以速度 \(v_0\)
射入一物块而不复射出，
求此后弹簧的最大压缩长度。

\paragraph{习题 4.12}
如图4.56，弹簧下悬挂质量分别为 \(m_1,m_2\)
的两个物体，开始时都处于静止状态。
突然把 \(m_1\) 与 \(m_2\) 的连线剪断后，
\(m_1\) 的最大速率是多大？
设弹簧的劲度系数
\(k=8.9\,\mathrm{N/m}\)，
\(m_1=500\,\mathrm g,\ m_2=300\,\mathrm g\)。

\paragraph{习题 4.13}
一质量为 \(m\) 的物体，从质量为 \(M\) 的圆弧形槽顶端
由静止滑下，该圆弧形槽半径为 \(R\)，
张角为 \(\pi/2\)（图4.57）。
如所有摩擦都可忽略，求：
（1）物体刚离开槽底端时，物体和槽的速度各是多少？
（2）在物体从A滑到B的过程中，物体对槽所做的功；
（3）物体到达B时对槽的压力。
\textit{图示：}槽放在无摩擦水平地面，圆弧从左上端A
滑至右下端B，高度差为 \(R\)，到B时切线水平。

\paragraph{习题 4.14}
证明：一个运动的小球与另一静止的质量相同的小球作弹性的非对心碰撞后，
它们总沿互成直角的方向离开。
（参考图4.24和图4.25。）

\paragraph{习题 4.15}
对于一维情况证明：
若两质点以某一相对速率靠近并做弹性碰撞，
那么碰撞后恒以同一相对速率离开，即
\[
 v_{10}-v_{20}=-(v_1-v_2).
\]

\paragraph{习题 4.16}
一质量为 \(m\) 的人造地球卫星沿一圆形轨道运动，
离开地面的高度等于地球半径的2倍（即 \(2R\)）。
试以 \(m,R\)、引力常量 \(G\)、地球质量 \(M\) 表示：
（1）卫星的动能；
（2）卫星在地球引力场中的引力势能；
（3）卫星的总机械能。

\paragraph{习题 4.17}
证明：行星在轨道上运动的总能量为
\[
 E=-\frac{GMm}{r_1+r_2},
\]
式中 \(M,m\) 分别为太阳和行星的质量，
\(r_1,r_2\) 分别为行星轨道的近日点和远日点距离。

\paragraph{习题 4.18}
发射地球同步卫星要利用“霍曼轨道”（图4.58）。
设发射一颗质量为 \(500\,\mathrm{kg}\) 的地球同步卫星，
先把它发射到高度为 \(1400\,\mathrm{km}\) 的停泊轨道上，
然后利用火箭推力使它沿此轨道的切线方向进入霍曼轨道。
霍曼轨道远地点即同步高度 \(36000\,\mathrm{km}\)，
在此高度上利用火箭推力使它进入同步轨道。
（1）先后两次火箭推力给予卫星的能量各是多少？
（2）先后两次推力使卫星的速率增加多少？
\textit{图示：}圆形近地停泊轨道→椭圆转移轨道→
圆形地球同步高度轨道。

\paragraph{习题 4.19}
两颗中子星质量都是 \(10^{30}\,\mathrm{kg}\)，
半径都是 \(20\,\mathrm{km}\)，相距 \(10^{10}\,\mathrm m\)。
如果它们最初都是静止的，试求：
（1）当它们的距离减少到一半时，
它们的速度各是多大？
（2）当它们就要碰上时，
它们的速度又将各是多大？

\paragraph{习题 4.20}
有一种说法认为地球上的一次灾难性物种（如恐龙）绝灭
是由于6500万年前一颗大的小行星撞入地球引起的。
设小行星半径 \(10\,\mathrm{km}\)，
密度为 \(6.0\times10^3\,\mathrm{kg/m^3}\)
（和地球的一样），
它撞入地球将释放多少引力势能？
这能量是唐山地震估计能量（见原书表4.1）的多少倍？
\textit{本书表4.1在教材PDF第131页给出1976年唐山大地震能量约\(1\times10^{18}\,\mathrm J\)。}

\paragraph{习题 4.21}
一个星体的逃逸速度为光速时，
意味着由于引力作用光子也不能从该星体表面逃离，
这星体就成了一个“黑洞”。
理论证明，对于这种情况，
逃逸速度公式 \(v_e=\sqrt{2GM/R}\) 仍然正确。
试计算太阳要成为黑洞，
它的半径应是多大（目前半径约为
\(R=7\times10^8\,\mathrm m\)）？
质量密度是多大？
比原子核的平均密度
\(2.3\times10^{17}\,\mathrm{kg/m^3}\) 大到多少倍？

\paragraph{习题 4.22}
理论物理学家霍金认为“黑洞”并不是完全“黑”的，
而是不断向外发射物质。
这种发射称为黑洞的“蒸发”。
他估计一个质量是太阳两倍的黑洞，
温度大约是 \(10^{-6}\,\mathrm K\)，
完全蒸发掉需要 \(10^{67}\) 年。
又据信宇宙大爆炸开始曾产生过许多微型黑洞，
但到如今这些微型黑洞都已由于蒸发而消失了。
若一个黑洞蒸发完所需时间和它的质量的3次方成正比，
而宇宙大爆炸发生在200亿年以前，
那么当时产生而至今已蒸发完的最大微型黑洞
的质量和半径各是多少？

\paragraph{习题 4.23}
\({}^{238}\mathrm U\) 核放射性衰变放出
\(\alpha\) 粒子时释放的总能量是 \(4.27\,\mathrm{MeV}\)，
求一个静止的 \({}^{238}\mathrm U\) 核放出的
\(\alpha\) 粒子的动能。

\paragraph{习题 4.24}
证明：把两体问题化为单体问题后，
一质点在另一质点参考系中的动能
等于两质点的内动能。

\paragraph{习题 4.25}
水平光滑桌面上放有质量分别为 \(M\) 和 \(m\)
的两个物体，二者用一根劲度系数为 \(k\)
的弹簧相连而处于静止状态。
今用棒击质量为 \(m\) 的物体，
使之获得一个指向另一物体的速度 \(v_0\)。
试利用约化质量概念求出此后弹簧的最大压缩长度。

\paragraph{习题 4.26}
在实验室内观察到相距很远的一个质子
（质量为 \(m_p\)）和一个氦核
（质量 \(M=4m_p\)）相向运动，
速率都是 \(v_0\)。
求二者能达到的最近距离
（忽略质子和氦核间的引力势能，
但二者间的电势能需计入。
电势能公式可根据引力势能公式猜出）。

\paragraph{习题 4.27}
有的黄河区段的河底高于堤外田地。
为了用河水灌溉堤外田地就用虹吸管越过堤面把河水引入田中。
虹吸管如图4.59所示，是倒U形，
其两端分别处于河内和堤外的水渠上。
如果河水水面和堤外管口的高度差是
\(5.0\,\mathrm m\)，
而虹吸吸管半径是 \(0.20\,\mathrm m\)，
则每小时可引入田地的河水体积是多少？
\textit{使用理想稳态不可压缩流体、忽略损耗的模型。}

\paragraph{习题 4.28}
喷药车的加压罐内杀虫药水表面的压强为
\(p_0=21\,\mathrm{atm}\)，
管道另一端喷嘴的直径为 \(0.8\,\mathrm{cm}\)
（图4.60）。
求喷药时每分钟喷出的杀虫药水的体积。
设喷嘴和罐内液面处于同一高度。
\textit{压强是绝对压强还是表压会改变数值，
应按题设 \(p_0\) 与大气出口压强的差计算。}
